% Part: methods
% Chapter: induction
% Section: structural-induction

\documentclass[../../../include/open-logic-section]{subfiles}

\begin{document}

\olfileid{mth}{ind}{sti}

\olsection{استقرای ساختاری}

تا اینجا از استقرا برای اثبات نتایج دربارهٔ همهٔ اعداد طبیعی استفاده
کرده‌ایم. اما می‌توان اصلی متناظر را مستقیماً برای اثبات نتایج دربارهٔ
همهٔ !!{element}های مجموعه‌ای که به‌طور استقرایی تعریف شده است به کار
برد. این اصل را اغلب \emph{استقرای ساختاری} می‌نامند، زیرا بر ساختار
اشیای به‌طور استقرایی تعریف‌شده تکیه دارد.

به‌طور کلی، تعریف استقرایی با (الف)~فهرستی از !!{element}های «آغازی»
مجموعه و (ب)~فهرستی از عملیات‌هایی داده می‌شود که از اعضای پیشین
مجموعه، !!{element}های تازه‌ای تولید می‌کنند. برای نمونه، در
مورد ترم‌های خوش‌ساخت، اشیای آغازی همان حروف‌اند. فقط یک عمل داریم:
\begin{align*}
  o(s_1, s_2) &= [s_1 \circ s_2]
\end{align*}
حتی می‌توانید خود اعداد طبیعی~$\Nat$ را حاصل یک تعریف استقرایی بدانید:
شیء آغازی~$0$ است و عمل همان تابع جانشین~$x + 1$ است.

برای اثبات چیزی دربارهٔ همهٔ اعضای یک مجموعهٔ به‌طور استقرایی
تعریف‌شده، یعنی برای اثبات اینکه هر !!{element} از مجموعه خاصیت~$P$ را
دارد، باید:
\begin{enumerate}
\item اثبات کنیم اشیای آغازی خاصیت~$P$ را دارند؛
\item اثبات کنیم برای هر عمل~$o$، اگر ورودی‌ها خاصیت~$P$ را داشته
  باشند، خروجی نیز آن را دارد.
\end{enumerate}
برای نمونه، برای اثبات چیزی دربارهٔ همهٔ ترم‌های خوش‌ساخت باید اثبات
کنیم آن چیز دربارهٔ همهٔ حروف صادق است، و اگر دربارهٔ هر یک از $s_1$
و $s_2$ صادق باشد، دربارهٔ $[s_1 \circ s_2]$ نیز صادق است.

\begin{prop}
  در هر ترم خوش‌ساخت~$t$، تعداد $[$ها با تعداد $]$ها برابر است.
\end{prop}

\begin{proof}
استقرای ساختاری به کار می‌بریم. ترم‌های خوش‌ساخت به‌طور استقرایی تعریف
شده‌اند: حروف اشیای آغازی‌اند و عمل~$o$ ترم‌های خوش‌ساخت تازه را از
ترم‌های پیشین می‌سازد.
\begin{enumerate}
\item گزاره برای هر حرف صادق است، زیرا تعداد $[$ها در یک حرف به‌تنهایی
  برابر~$0$ است و تعداد $]$ها در آن نیز~$0$ است.
\item فرض کنید تعداد $]$ها در $s_1$ با تعداد $[$ها در $s_1$ برابر
  باشد، و همین امر دربارهٔ $s_2$ نیز صادق باشد. تعداد $[$ها در
  $o(s_1, s_2)$، یعنی در $[s_1 \circ s_2]$، برابر است با مجموع تعداد
  $[$ها در $s_1$ و $s_2$، به‌اضافهٔ یک. تعداد $]$ها در $o(s_1, s_2)$
  نیز برابر است با مجموع تعداد $]$ها در $s_1$ و $s_2$، به‌اضافهٔ یک.
  پس تعداد $[$ها در $o(s_1, s_2)$ با تعداد $]$ها در $o(s_1,s_2)$
  برابر است.
\end{enumerate}
\end{proof}

\begin{prob}
  با استقرای ساختاری اثبات کنید هیچ ترم خوش‌ساختی با~$]$ آغاز نمی‌شود.
\end{prob}

اثبات دیگری با استقرای ساختاری ارائه کنیم. یک پاره‌رشتهٔ آغازی سره از
رشتهٔ نمادهای~$t$، هر رشتهٔ~$s$ای است که هنگام خواندن از چپ، نماد به
نماد با $t$ یکسان است، اما $t$ بلندتر است. برای نمونه،
$[a \circ {}$ یک پاره‌رشتهٔ آغازی سره از $[a \circ b]$ است؛ اما
$[b \circ {}$ چنین نیست، چون آن دو در نماد دوم تفاوت دارند، و
$[a \circ b]$ نیز چنین نیست، چون طولشان برابر است.

\begin{prop}\ollabel{prop:initial}
  در هر پاره‌رشتهٔ آغازی سره از یک ترم خوش‌ساخت~$t$، تعداد $[$ها از
  تعداد $]$ها بیشتر است.
\end{prop}

\begin{proof}
  بر~$t$ استقرا می‌کنیم:
  \begin{enumerate}
  \item $t$ یک حرف است: آنگاه $t$ هیچ پاره‌رشتهٔ آغازی سره‌ای ندارد.
  \item برای دو ترم خوش‌ساخت $s_1$ و~$s_2$،
    $t = [s_1 \circ s_2]$. اگر $r$ یک پاره‌رشتهٔ آغازی سره از $t$
    باشد، چند حالت ممکن است:
    \begin{enumerate}
    \item $r$ فقط $[$ است: آنگاه $r$ یک $[$ بیشتر از $]$ دارد.
    \item $r$ برابر $[r_1$ است، که در آن $r_1$ یک پاره‌رشتهٔ آغازی
      سره از~$s_1$ است: چون $s_1$ ترمی خوش‌ساخت است، بنا بر فرض
      استقرا $r_1$ تعداد $[$ بیشتری از $]$ دارد، و همین امر برای
      $[r_1$ صادق است.
    \item $r$ برابر $[s_1$ یا $[s_1 \circ {}$ است: بنا بر نتیجهٔ
      پیشین، تعداد $[$ها و $]$ها در~$s_1$ برابر است؛ پس تعداد $[$ها در
      $[s_1$ یا $[s_1 \circ {}$ یک واحد بیشتر از تعداد~$]$هاست.
    \item $r$ برابر $[s_1 \circ r_2$ است، که در آن $r_2$ یک پاره‌رشتهٔ
      آغازی سره از~$s_2$ است: بنا بر فرض استقرا، $r_2$ تعداد $[$ بیشتری
      از $]$ دارد. بنا بر نتیجهٔ پیشین، تعداد $[$ها و $]$ها در~$s_1$
      برابر است. پس در $[s_1 \circ r_2$ تعداد $[$ها از تعداد~$]$ها
      بیشتر است.
    \item $r$ برابر $[s_1 \circ s_2$ است: بنا بر نتیجهٔ پیشین، تعداد
      $[$ها و $]$ها در $s_1$ برابر است و همین امر برای~$s_2$ نیز صادق
      است. پس در $[s_1 \circ s_2$ یک $[$ بیشتر از $]$ وجود دارد.
    \end{enumerate}
  \end{enumerate}
\end{proof}

\end{document}
